\zlabel{firstpagetocount}       % DO NOT REMOVE! Used for counting number of pages of main text

% Part Labels Explanation:
% Labels in the form \label{chapter:<part_name>} are used to categorize chapters into specific parts of the thesis structure.
% These labels are essential for calculating the ratio of text dedicated to different thesis sections (e.g., Introduction, Methodology, Results, Discussion, Summary).
% Allowed labels:
% - chapter:introduction -> Introduction (<10% of the text).
% - chapter:method       -> Methodology (<20% of the text).
% - chapter:results      -> Results (30–40% of the text).
% - chapter:discussion   -> Analysis, Discussion, and Conclusions (30–40% of the text).
% - chapter:summary      -> Summary (<0.5 pages).
% How to use: Add a \label{chapter:<part_name>} for each chapter to indicate its part. 
% These labels ensure consistency and allow automated tests to validate the thesis structure.

\chapter{Sissejuhatus}\label{chapter:introduction} % This command creates a numbered chapter titled "Sissejuhatus" (Estonian for "Introduction"). The \label{chapter:introduction} assigns a label to the chapter, allowing you to reference it later (e.g., \ref{chapter:introduction}).
Euroopa liit on endale seadnud eesmärgiks aastaks 2050 saavutada kliimaneutraalsus.
Selle eesmärgi täitmiseks on Euroopa Liidus loodud erinevaid rahastusvehendeid.
Üheks selliseks vahendiks on Euroopa ühendamise rahastu (CEF), mis investeerib kolme 
peamisesse suunda: energeetika, transport ja digivõrgud. \cite{cef}

CEF transport algatas 2023 aastal eFTI4EU projekti,
mille eesmärgiks on standardiseeritult digitaliseerida veoinfo vahetus pädevate 
asutuste ning vedajate vahel lähtudes Euroopa Parlamendi määrusest 2020/1056 \cite{efti4eu_project}. 

eFTI4EU jätkuks loodi 2024 aastal eFTI4ALL projekt, 
mis kasutab ning täiendab eelnevas projektis loodud lahendusi päriselulistes olukordades. 
\cite{efti4all_project}

\section{Taust}

eFTI4EU projekti üks suuremaid tulemusi oli eFTI näidislahenduse loomine, 
mis sisaldab kolme põhilist eFTI arhitektuuri komponenti - eFTI värav, eFTI näidisplatvorm 
ja pädevate asutuste juurdepääsupunkt. \cite{efti4eu_reference_impl}

eFTI värav on süsteemi keskne osa, millega suhtlevad nii eFTI platvormid kui ka pädevad asutused.
Väravaid haldavad kõik liikmesriigid ise, ning on kohustatud pakkuma seda avaliku 
teenusena transpordiettevõtetele ja pädevatele asutustele. Värava peamine roll on andmete ning
päringute edastamine osapoolte vahel.Väravad peavad olema võimelised omavahel päringuid vastu võtma.
Selleks, et standardiseerida suhtlus väravate vahel on kasutusele võetud eDelivery andmevahetuslahendus.
eDelivery on avatud spetsifikatsioonidel põhinev infrastruktuur ja protokoll, 
mis kasutab AS4 (Applicability Statement 4) standardit sõnumite vahetamiseks. 
See on nagu digitaalne postisüsteem, mis tagab, et sõnumid liiguvad ühest punktist
teise kontrollitud ja krüpteeritud viisil. \cite{ria_edelivery_selgitus}

Platvorm täidab eFTIs andmete looja ja haldaja rolli. Platvormi kaudu on võimalik luua eFTI
andmehulki, mis kirjeldavad veost ning selle liikumist. Platvormid on üldjuhul eraettevõtted,
kes soovivad muuta oma veoandmed kättesaadavaks pädevatele asutustele läbi eFTI värava. 
Värav saab pärida ainult alamhulka eFTI andmetest, mis sisaldab vaid neid andmeid, mida pädev asutus
seaduslikult peaks kätte saama. Seeläbi minimeeritakse liigsete andmete edastamist.
eFTI andmete alamhulgad on defineerutud Euroopa Liidu määrusega 2024/2024 \cite{eu_reg_2024_2024}

Pädevate asutuste juurdepääsupunktid ehk Authority Access Points (AAP) võimaldavad huvitatud osapooltel pärida eFTI andmeid.
Vastavalt liikmesriigi vajadustele võib AAP liidestus väravaga olla mistahes sugune.
Eesti kontekstis on võimalik kasutada X-Tee sõnumeid, et pärida andmeid eFTI keskkonnast. \cite{dlk_2022_efti_analysis}
Kuna AAP-il pole otsest kokkupuudet eDelivery-ga siis selle lõputöö raamidest jääb AAP välja.

\section{Probleem}

eFTI4EU konsortiumi poolt loodud näidisarhitektuuris on mõningaid puudujääke. 
Nende puuduste hulgast selle lõputöö kontekstis kõige olulisem on eDelivery SMP/SML kasutamine platvormi ning värava vahel.
Tänane lahendus võimaldab eDelivery põhist suhtlust värava ja platvormi vahel kasutades staatilist konfiguratsiooni,
mis ei vasta euroopa liidu regulatsioonile 2024/1942 artikkel 9 (3) \cite{eu_reg_2024_1942}.
Hetkene lahendus töötab sarnaselt sellele kuidas väravad omavahel suhtlevad - staatiliselt.







 %  This command inserts the content from the file introduction.tex located in the chapters folder.

\zlabel{lastpagetocount}        % DO NOT REMOVE! Used for counting number of pages of main text
